\section{Canais e alinhamento com a proposta de valor}
\label{sec:canais}

O PetCuida é uma rede local de cuidado preventivo que conecta tutores, serviços comunitários e prestadores veterinários. Sua proposta de valor é \textit{cuidar antes, pagar de forma possível e encontrar orientação confiável}. O desenho dos canais, portanto, precisa reduzir barreiras financeiras e informacionais, alcançar tutores de baixa renda e sustentar a continuidade do cuidado.

\subsection{Mapeamento de canais}
\label{subsec:mapeamento}

\begin{longtable}{p{0.25\textwidth}p{0.67\textwidth}}
\toprule
\textbf{Categoria} & \textbf{Canais possíveis para o PetCuida}\\
\midrule
\endhead
Comunicação & Instagram, Facebook e WhatsApp; grupos de bairro; parcerias com ONGs, pet shops e clínicas; eventos de adoção e feiras comunitárias; indicação boca a boca; cartazes em pontos de circulação; conteúdo educativo sobre prevenção; imprensa e influenciadores locais.\\
\addlinespace
Distribuição/entrega & Aplicativo/webapp PetCuida; WhatsApp para cadastro e lembretes; clínicas e consultórios parceiros; pontos de apoio de ONGs; agenda de serviços comunitários; plataformas de pagamento híbrido (créditos + dinheiro).\\
\addlinespace
Relacionamento & Atendimento humano por WhatsApp; chatbot para dúvidas e lembretes; suporte pós-atendimento; notificações de vacinação e vermifugação; comunidade online; pesquisas de satisfação; acompanhamento de retornos e histórico digital do pet.\\
\bottomrule
\end{longtable}

\subsection{Estratégia: canais prioritários}
\label{subsec:estrategia}

A priorização considera três critérios: alinhamento com a proposta de valor, adequação ao cliente-alvo e viabilidade no estágio inicial da startup. Os nove canais abaixo formam um primeiro portfólio enxuto, mensurável e compatível com uma operação local.

\begin{table}[htbp]
\centering\small
\caption{Canais prioritários por categoria}
\label{tab:prioritarios}
\begin{tabularx}{\textwidth}{p{0.19\textwidth}p{0.25\textwidth}Y}
\toprule
\textbf{Categoria} & \textbf{Canal prioritário} & \textbf{Papel no funil}\\
\midrule
Comunicação & WhatsApp e grupos de bairro & Alcance direto, linguagem acessível e circulação de indicações.\\
Comunicação & Parcerias com ONGs e clínicas & Empresta confiança local e aproxima o público já atendido por organizações do ecossistema.\\
Comunicação & Instagram/Facebook com conteúdo educativo & Gera descoberta, esclarece prevenção e permite segmentação geográfica de baixo custo.\\
\addlinespace
Distribuição & Aplicativo/webapp PetCuida & Centraliza calendário, triagem orientativa, prontuário e localização de parceiros.\\
Distribuição & WhatsApp para onboarding e lembretes & Reduz barreira tecnológica e entrega lembretes no canal já usado pelo tutor.\\
Distribuição & Clínicas veterinárias parceiras & Converte a intenção em atendimento e viabiliza preços sociais e créditos.\\
\addlinespace
Relacionamento & Atendimento humano por WhatsApp & Oferece acolhimento e encaminhamento em casos que exigem contexto.\\
Relacionamento & Notificações de prevenção & Mantém o cuidado ativo com lembretes de vacinação, vermifugação e retorno.\\
Relacionamento & Comunidade com ONGs e tutores & Sustenta confiança, educação e recorrência por meio de apoio local.\\
\bottomrule
\end{tabularx}
\end{table}
